\section{Seguridad}

\begin{itemize}[leftmargin=*,itemsep=4pt]
  \item \textbf{AgroDash es sólo lectura, y se fuerza.} Toda conexión fija
    \cd{SET SESSION CHARACTERISTICS AS TRANSACTION READ ONLY} y sólo se ejecutan
    \cd{SELECT}. No se confía en que el rol esté bien configurado del otro lado.

  \item \textbf{RLS en modo cierre} en las tablas públicas de Supabase, deliberadamente
    \emph{sin políticas}: sólo los roles de servicio acceden y la API REST pública queda
    bloqueada. Se verificó que no rompe al forecaster, que corre con un rol que la
    salta.

  \item \textbf{Vistas con \cd{security\_invoker}}, para que ejecuten con los permisos de
    quien consulta y no con los del dueño.

  \item \textbf{Claves sólo por entorno}, nunca en el código ni en el repositorio.
    Ningún secreto está versionado.

  \item \textbf{Comparación en tiempo constante} de las API keys
    (\cd{secrets.compare\_digest}), para no filtrar la clave por diferencias de tiempo de
    respuesta.

  \item \textbf{Clave exigida sólo si está configurada.} En local los servicios corren
    abiertos para no estorbar; en producción la variable está definida y el borde exige
    el encabezado. La excepción son \cd{/health} y \cd{/salud/ingesta}, abiertas a
    propósito para monitoreo automático: no exponen datos de la serie.

  \item \textbf{Los servicios no se publican; se publica el borde.} Los dos agentes y la
    réplica escuchan en \cd{127.0.0.1} y nada más. Lo único alcanzable desde internet es
    nginx, con TLS. Hasta la mudanza los servicios escuchaban en \cd{0.0.0.0} y lo único
    que los tapaba era el grupo de seguridad de la nube: cualquier máquina de la misma
    red privada los alcanzaba. Atarlos al loopback quita esa clase entera de error.

  \item \textbf{Freno de fuerza bruta en el borde.} El límite de intentos de la propia
    consola vive en memoria del proceso, así que un reinicio lo borra y en un despliegue
    con varias instancias no se comparte. La defensa que no depende de eso es doble: cada
    intento cuesta un derivado de clave deliberadamente caro (unos 400\,ms de CPU), y
    nginx limita la ruta de acceso a seis intentos por minuto. Verificado, no sólo
    configurado: diez intentos seguidos dan seis 401 y cuatro 429.

  \item \textbf{Falla cerrada.} Sin contraseña configurada, la consola en producción no
    se abre: responde 503 y lo dice. Un despliegue anterior quedó expuesto justamente por
    hacer lo contrario, abrirse cuando faltaba la variable.

  \item \textbf{Errores internos genéricos.} Un fallo interno responde el 500 estándar;
    el detalle queda en el log del servidor, nunca en la respuesta.
\end{itemize}

\subsection{Secretos necesarios}

\begin{center}
\small
\begin{tabular}{@{}p{44mm}p{96mm}@{}}
\toprule
\textbf{Variable} & \textbf{Para qué} \\
\midrule
\cd{DATABASE\_URL} & Conexión a Supabase (usar el \emph{Session pooler}; la conexión
  directa es sólo IPv6). \\
\cd{STORE\_URL} & El store del agente, separado de la fuente a propósito. \\
\cd{ANTHROPIC\_API\_KEY} & Los dos agentes. \\
\cd{FORECAST\_API\_KEY} & Proteger el servicio de pronóstico. \\
\cd{HISTORICO\_API\_KEY} & Proteger el servicio del Histórico. \\
\cd{DEBUGGER\_PASSWORD} & Entrar a la consola. \\
\cd{DEBUGGER\_SESSION\_SECRET} & Firmar la cookie de sesión. \\
\bottomrule
\end{tabular}
\end{center}

\subsection{Pendiente de seguridad}

Rotar la clave débil del rol de sólo lectura que se usó para probar el acceso a la base
viva de Cartago por la malla Tailscale.
